\section{Билет 61. КПД тепловых машин. Цикл Карно. Вычисление КПД цикла Карно термодинамически. Теорема Карно.}

\begin{definition}
\textbf{Коэффициентом полезного действия (КПД)} тепловой машины называется отношение полезной работы $A$, совершаемой за цикл, к количеству теплоты $Q_1$, полученному от нагревателя:
\begin{equation}
    \eta = \frac{A}{Q_1} = \frac{Q_1 - |Q_2|}{Q_1} = 1 - \frac{|Q_2|}{Q_1}, \label{eq:efficiency_61}
\end{equation}
где $|Q_2|$ --- количество теплоты, отданное холодильнику.
\end{definition}

\begin{definition}
\textbf{Циклом Карно} называется идеальный термодинамический цикл, состоящий из двух изотермических и двух адиабатических процессов, совершаемых обратимо (квазистатически). Цикл работает между двумя тепловыми резервуарами с постоянными температурами: нагревателя $T_1$ и холодильника $T_2$ ($T_1 > T_2$).
\begin{itemize}
    \item 1--2: Изотермическое расширение при $T_1$ (получение теплоты $Q_1$);
    \item 2--3: Адиабатическое расширение (температура падает от $T_1$ до $T_2$);
    \item 3--4: Изотермическое сжатие при $T_2$ (отдача теплоты $|Q_2|$);
    \item 4--1: Адиабатическое сжатие (температура повышается от $T_2$ до $T_1$).
\end{itemize}
\end{definition}

\begin{theorem}[Вычисление КПД цикла Карно]
Для идеального газа количество теплоты в изотермических процессах равно работе:
\begin{align}
    Q_1 &= \nu R T_1 \ln\frac{V_2}{V_1}, \\
    |Q_2| &= \nu R T_2 \ln\frac{V_3}{V_4}.
\end{align}
Для адиабатических процессов 2--3 и 4--1 справедливы соотношения $TV^{\gamma-1} = \text{const}$:
\begin{equation}
    T_1 V_2^{\gamma-1} = T_2 V_3^{\gamma-1}, \quad T_1 V_1^{\gamma-1} = T_2 V_4^{\gamma-1}.
\end{equation}
Разделив первое равенство на второе, получаем:
\begin{equation}
    \left(\frac{V_2}{V_1}\right)^{\gamma-1} = \left(\frac{V_3}{V_4}\right)^{\gamma-1} \quad \Rightarrow \quad \frac{V_2}{V_1} = \frac{V_3}{V_4}.
\end{equation}
Следовательно, логарифмы в выражениях для $Q_1$ и $|Q_2|$ равны. Подставляя их в формулу КПД \eqref{eq:efficiency_61}:
\begin{equation}
    \eta_{\text{К}} = 1 - \frac{\nu R T_2 \ln(V_2/V_1)}{\nu R T_1 \ln(V_2/V_1)} = 1 - \frac{T_2}{T_1}. \label{eq:carnot_eff_61}
\end{equation}
\end{theorem}

\begin{theorem}[Теорема Карно]
\begin{enumerate}
    \item КПД любой тепловой машины, работающей между нагревателем с температурой $T_1$ и холодильником с температурой $T_2$, не может превышать КПД идеальной машины Карно: $\eta \le \eta_{\text{К}}$.
    \item КПД всех обратимых машин, работающих между одинаковыми температурами $T_1$ и $T_2$, одинаков и равен $\eta_{\text{К}} = 1 - T_2/T_1$, независимо от природы рабочего тела и конструкции машины.
    \item КПД необратимых (реальных) машин всегда строго меньше $\eta_{\text{К}}$.
\end{enumerate}
\end{theorem}

\begin{remark}
\textbf{Физические следствия:}
\begin{itemize}
    \item $\eta_{\text{К}} = 1$ возможно только при $T_2 = 0$ К (абсолютный ноль недостижим по третьему началу термодинамики) или $T_1 \to \infty$ (практически невозможно).
    \item Для повышения КПД реальных двигателей стремятся максимально повысить температуру нагревателя $T_1$ (жаропрочные сплавы, турбокомпрессоры) и понизить температуру холодильника $T_2$ (улучшенное охлаждение).
    \item Цикл Карно является теоретическим пределом эффективности преобразования теплоты в работу и служит эталоном для оценки реальных термодинамических циклов.
\end{itemize}
\end{remark}
